\section{References}

\subsection*{Sacred Scripture and Ancient Versions}

\begin{itemize}
\item \textbf{Douay-Rheims Bible (Challoner revision).} \emph{The Holy Bible, Translated from the Latin Vulgate} (1749--1752; public domain). Primary English Scripture text throughout: Genesis 1:5; 2:21; 13:14--17; 15:1--21; 22:17; 26:4; Exodus 32:13; Deuteronomy 1:10; 10:22; 28:62; 1 Paralipomenon 27:23; 2 Esdras (Nehemiah) 9:23; Romans 4:3, 18--22; 1 Corinthians 15:41; Galatians 3:6--9; Philippians 2:15; Hebrews 11:1, 11--12.
\item \textbf{Clementine Vulgate.} \emph{Biblia Sacra Vulgata}, Sixto-Clementine edition (public domain). Latin text of Genesis 15:1--21.
\item \textbf{Kittel, Rudolf, and Georg Beer, eds.} \emph{Biblia Hebraica}, Pars I. Leipzig: J. C. Hinrichs, 1909 (public domain). Genesis 15:5, 12, 17--18, pp.~18--19.
\item \textbf{Swete, Henry Barclay, ed.} \emph{The Old Testament in Greek according to the Septuagint}, vol.~I: Genesis--IV Kings, 4th ed. Cambridge: University Press, 1909 (public domain). Genesis 15:1--21 (supplying Codex Alexandrinus across Codex Vaticanus's opening lacuna Gen 1:1--46:28a).
\item \textbf{Brenton, Lancelot Charles Lee, trans.} \emph{The Septuagint Version of the Old Testament, According to the Vatican Text}. London: Samuel Bagster and Sons, 1844 (public domain). Genesis 15:1--21.
\item \textbf{Syriac Peshitta.} \emph{Peshitta Old Testament: Genesis}, 15:5, 12, 17--18 (ancient Syriac text, public domain; verified via recorded digital-delivery passage record, printed exemplar unnamed).
\end{itemize}

\subsection*{Second Temple and Jewish Sources}

\begin{itemize}
\item \textbf{Book of Jubilees.} Translated by R. H. Charles, in \emph{The Apocrypha and Pseudepigrapha of the Old Testament}, vol.~II. Oxford: Clarendon Press, 1913 (public domain; verified via recorded passage record). Chapter 14:1--20.
\item \textbf{Josephus, Flavius.} \emph{Antiquitates Judaicae} (\emph{Antiquities of the Jews}), Book I, ch.~10, \S 3. Translated by William Whiston (London, 1737; public domain; verified via recorded passage record).
\item \textbf{Abarbanel, Isaac.} \emph{Perush 'al ha-Torah} on Bereshit (Genesis) 15, opening formula and prefatory Questions 7 and 16. Warsaw, 1862 (public domain; directly inspected checked Hebrew text).
\end{itemize}

\subsection*{Patristic, Medieval, and Catholic Authors}

\begin{itemize}
\item \textbf{Clement of Rome, St.} \emph{First Epistle to the Corinthians}, ch.~10. In \emph{The Ante-Nicene Fathers}, vol.~1, eds. Alexander Roberts and James Donaldson. Buffalo: Christian Literature Publishing Co., 1885 (public domain).
\item \textbf{Ambrose of Milan, St.} \emph{De excessu fratris Satyri}, Book II (\emph{On the Belief in the Resurrection}), \S 95. Translated by H. de Romestin, in \emph{Nicene and Post-Nicene Fathers}, Series 2, vol.~10. Edinburgh: T. \& T. Clark, 1896 (public domain).
\item \textbf{Jerome, St.} \emph{Quaestiones Hebraicae in Genesim}, ch.~15. Edited by Domenico Vallarsi, in J.-P. Migne, \emph{Patrologia Latina}, vol.~23 (public domain).
\item \textbf{Augustine of Hippo, St.} \emph{De Civitate Dei}, Books X.8 and XVI.23--24. Latin text edited by Emanuel Hoffmann, \emph{Corpus Scriptorum Ecclesiasticorum Latinorum} (CSEL), vol.~40, parts 1--2 (Vienna: F. Tempsky, 1899--1900; public domain); English translation by Marcus Dods, \emph{The City of God}, 2 vols. (Edinburgh: T. \& T. Clark, 1871; public domain).
\item \textbf{Augustine of Hippo, St.} \emph{De Genesi ad litteram}, Book XII, ch.~6. In J.-P. Migne, \emph{Patrologia Latina}, vol.~34 (public domain).
\item \textbf{John Chrysostom, St.} \emph{Homilies on the Epistle to the Romans}, Homily 9 (on Rom 4:18--20), and \emph{Commentary on the Epistle to the Galatians}, on Gal 3:6. In \emph{Nicene and Post-Nicene Fathers}, Series 1, vols.~11 and 13, ed. Philip Schaff. Edinburgh: T. \& T. Clark (public domain).
\item \textbf{Bede the Venerable, St.} \emph{In principium Genesis}, on Genesis 15:1--17. In J.-P. Migne, \emph{Patrologia Latina}, vol.~91 (public domain).
\item \textbf{Alcuin of York.} \emph{Interrogationes et responsiones in Genesim}, Questions 165--170. In J.-P. Migne, \emph{Patrologia Latina}, vol.~100 (public domain).
\item \textbf{Angelomus of Luxeuil.} \emph{Commentarius in Genesim}, on ch.~15. In J.-P. Migne, \emph{Patrologia Latina}, vol.~115 (public domain).
\item \textbf{Thomas Aquinas, St.} \emph{Summa Theologiae}, II-II, q.~173, a.~3; II-II, q.~174, a.~6; III, q.~30, a.~3, ad~3. Translated by the Fathers of the English Dominican Province (London, 1920; public domain).
\item \textbf{Thomas Aquinas, St.} \emph{In omnes D. Pauli Apostoli epistolas commentaria}, vol.~2, on Galatians 3:16. Mechlin: H. Dessain, 1857 (public domain).
\item \textbf{Lapide, Cornelius a, S.J.} \emph{Commentaria in Pentateuchum Mosis}. Antwerp, 1700 (public domain; directly inspected checked Latin text of notes on Genesis 15:1, 5, 11, 12, and 17).
\end{itemize}
